% Include this section with graphicx, amsmath, booktabs, float, and natbib available.
% SURDResultsPath is the path to this folder from the compiling document.
\providecommand{\SURDResultsPath}{collaborator_rebuild/results_section}

\section{Campaign resolved SURD components and sampling sensitivity}
\label{sec:campaign_surd_results}

\subsection{Observed series and analysis scope}

We analyze historical observations of NGC~5548 in five observing campaigns
from 1989 through 1993. The adopted line sample contains 242 spectra on
224 distinct dates. Corrected published total, blue, core, and red
H$\beta$ measurements are used directly from the table accompanying
\citet{CRWanders1996}, together with 557 original optical continuum
measurements at 5100~\AA. Duplicate spectral dates are represented by equal
means, with mean uncertainties obtained from the quadrature sum of quoted
errors divided by the number of measurements. Shared calibration covariance
is unavailable. The original continuum retains host starlight.
All four line variables share spectra, and the total nearly equals the sum
of the three components. We retain the published rounding and do not treat
these measurements as independent observations of four physical regions.

Campaign membership follows the intervals in Table~2 of
\citet{CRPeterson2004}: JD less 2400000 from 47509 through 47809,
47861 through 48179, 48225 through 48534, 48623 through 48898,
and 48954 through 49255. The blue, core, and red definitions follow
the observed wavelength intervals in source Section~3.2, respectively
4870 through 4920, 4920 through 4960, and 4960 through 5010~\AA.
The source figure caption gives a different blue lower boundary;
using the published measurements avoids claiming an exact reconstruction
of the integrations from the archived profile pixels.

For each campaign, we analyze two variable sets: total, blue, core, and
red H$\beta$; and continuum, blue, core, and red H$\beta$.
Every variable in a set becomes the future target in turn, while all four
past variables, including the target's own history, remain predictors.
We use the unchanged SURD implementation of \citet{CRSURD2024}, with
the averaging and normalization described by \citet{CRThesis2025}.
For four predictors, the listed combinations give four unique,
eleven redundant, and eleven synergistic components, hence 26 in total.

For future target $Y(t)$ and predictor vector $\mathbf{X}_{\tau}(t)$,
we write each normalized component as
\begin{equation}
 f_a(\tau)=\frac{I_a(Y;\mathbf{X}_{\tau})}
 {I(Y;\mathbf{X}_{\tau})},
 \qquad \sum_{a=1}^{26}f_a(\tau)=1,
 \label{eq:cr_component_normalization}
\end{equation}
when joint mutual information is nonzero. Information leakage uses a
separate denominator,
\begin{equation}
 \ell(\tau)=\frac{H(Y\mid\mathbf{X}_{\tau})}{H(Y)}.
 \label{eq:cr_leakage}
\end{equation}
The component fractions therefore describe the allocation of available
predictive information, while leakage describes the fraction of target
entropy left unexplained. A large normalized component can coexist with
substantial leakage.

The saved exploratory scans cover base lags from zero through 30 observed
days. Predictors are interpolated at native target dates within the same
campaign, with no extrapolation and a maximum interpolation gap of ten days.
We retain a minimum of 32 finite tuples as a numerical guard.
Probability histograms use either two or three marginal quantile bins,
whose boundaries are determined from each full native campaign series.
The two additional scenarios place the blue or red predictor five days
earlier than the common base lag. These stipulated offsets are sensitivity
choices, not measured delays. Their construction gives one lag to each
predictor; it does not execute the expanded predictor vector of thesis
equation~3.18.

\subsection{Campaign support and sensitivity to bin choice}

Each bin setting contains 3720 configurations, with 1402 passing the
numerical guard. At positive lags there are 1326 supported configurations.
The native inputs, settings apart from bin count, tuple counts, and support
status agree between the two scans. This permits a direct comparison of
bin settings at each lag. It does not hold target dates fixed across lags.

We quantify the change in the normalized component distribution using
\begin{equation}
 D_{\rm TV}(\tau)=\frac{1}{2}\sum_{a=1}^{26}
 \left|f_a^{(2)}(\tau)-f_a^{(3)}(\tau)\right|.
 \label{eq:cr_tv}
\end{equation}
The superscripts denote bin count. Leakage changes are reported as
$|\ell^{(2)}-\ell^{(3)}|$. Both quantities are descriptive diagnostics;
neither is a significance test or an uncertainty estimate.

\begin{table}[H]
\centering
\caption{Native date counts and positive lag support by campaign.
Each campaign has 720 positive lag configurations per bin setting.
$N_{\rm line}$ is the number of native dates for each line series and
$N_F$ is the continuum count. $N_{\rm usable}$ applies to the earlier
support that changes with lag. The last column gives the range of retained
date counts after intersection across all positive lags in each curve.}
\label{tab:cr_campaign_support}
\small
\input{\SURDResultsPath/campaign_rows.tex}
\end{table}

Table~\ref{tab:cr_campaign_support} shows that numerical support differs
strongly between campaigns. Only one positive lag configuration is usable
in 1991. All 720 configurations in 1993 pass the guard, but the median
$D_{\rm TV}$ is 0.515 and every supported configuration in that campaign
has $D_{\rm TV}>0.2$. Across all campaigns, median $D_{\rm TV}$ is 0.484
and median absolute leakage change is 0.182. Thus even the campaign with
complete numerical coverage is sensitive to the probability discretization.
Complete coverage under the tuple guard does not establish reliable
estimation of the component fractions.

\subsection{Holding target dates fixed across all positive lags}

To remove changes in native target support along a lag curve, we retain
only target dates at which every required predictor is finite at every
base lag from one through 30 days. This intersection is computed separately
for each campaign, variable set, target, and scenario. The resulting dates
are used at every lag and for both bin settings. The lag range, interpolation
gap limit, numerical guard, and full campaign quantile boundaries are
unchanged. Zero lag is excluded from the intersection and comparison.

Only four of 120 curves retain at least 32 target dates: the common
continuum target in 1992, and the common, blue earlier, and red earlier
continuum targets in 1993. Every future line target fails the guard under
this stricter construction. The 116 excluded curves have no decomposition
in this run. The four retained curves produce 240 decompositions across
30 lags and two bin settings, with 6240 component rows.
Normalization is defined in all executed cases, with maximum absolute
component sum error $1.22\times10^{-15}$; this verifies arithmetic only.

\begin{table}[H]
\centering
\caption{Sensitivity between bin settings with identical target dates
across all positive lags. Every row has future continuum as the target.
Each median uses 30 base lags. The last column counts components whose
maximizing lag sets do not overlap between bin settings. All components,
including weak ones, are counted; ties are retained within absolute
tolerance $10^{-12}$.}
\label{tab:cr_fixed_support}
\small
\input{\SURDResultsPath/fixed_rows.tex}
\end{table}

The sensitivity persists after fixing the target dates
(Table~\ref{tab:cr_fixed_support}). Median $D_{\rm TV}$ ranges from
0.523 to 0.610. In addition, changing the retained observations affects the
estimates at a fixed bin count. For the 1993 common curve, median
$D_{\rm TV}$ between the earlier varying support estimates and the new fixed
support estimates is 0.475 for two bins and 0.391 for three bins.
Across the four retained curves and both bin settings, these medians range
from 0.272 to 0.475. This measures sensitivity to sample selection and can
include temporal variation; it does not isolate estimator bias.

The retained sample has 36 dates for the 1992 curve and 78 through 90
dates for the 1993 curves. In the 1993 common case, median occupied
histogram cell counts are 14.5 for two bins and 25 for three bins, with
90 tuples. The corresponding medians of singleton cell counts are four
and ten. Repeated or interpolated observations are not independent samples,
and structural dependence can also leave cells empty. These occupancy
counts are diagnostics rather than estimator validation.

\subsection{Individual components and descriptive extrema}

Figures~\ref{fig:cr_1992_common} through \ref{fig:cr_1993_red}
show all individual components and leakage for the four retained curves.
The notation uses $F$ for continuum, $B$ for blue, $C$ for core, and $R$
for red; the outer operator identifies a unique, redundant, or synergistic
component. Dotted curves use two bins and solid curves use three bins.
All component panels share the joint information denominator in
equation~\ref{eq:cr_component_normalization}, while the leakage panel uses
equation~\ref{eq:cr_leakage}.

The 1993 common case is selected as the example with the largest retained
date count, rather than by the amplitude or location of a peak.
Table~\ref{tab:cr_extrema} reports the largest component value within each
information type over the declared positive lag range.
The largest unique contribution comes from continuum history in both
bin settings, but its maximizing lag changes from six days to one day.
The largest redundancy changes from $R(B,C)$ at 25 days to $R(B,C,R)$
at eleven days. The largest synergy changes from $S(F,R)$ at the
30 day scan boundary to $S(F,C)$ at four days. The boundary maximum
cannot identify an interior characteristic delay. Component identities,
amplitudes, and maximizing lags therefore change with discretization
despite identical native target support.

\begin{table}[H]
\centering
\caption{Descriptive maxima in the 1993 common scenario with 90 fixed
continuum target dates. Each row maximizes over all components of the
specified type and all base lags from one through 30 days. These maxima
have no null calibration or delay uncertainty and are not physical lag
measurements.}
\label{tab:cr_extrema}
\small
\input{\SURDResultsPath/extrema_rows.tex}
\end{table}

\subsection{Ideal categorical estimator calibration}

We check the same four predictor estimator against eight exact categorical
laws: independent variables, a target equal to one predictor, a target
duplicated by two predictors, and a target given by the modular sum of two
independent predictors, each with two or three states. The dependent laws
have known unique, redundant, or synergistic information, respectively.
Their exact joint information is the base two logarithm of state count
and their leakage is zero. The independent law has zero joint information
and unit leakage, so its exact normalized component fractions are undefined.
The official routine reproduces these exact reference values in the
declared predictor order.

We draw 200 independent samples per law and sample size with fixed seed
20261007, giving 6400 decompositions. Sizes of 36, 78, and 90 match
the retained curve counts; 512 is an artificial reference.
Table~\ref{tab:cr_calibration} shows substantial apparent information under
independence, particularly with three states. This benchmark demonstrates
finite sample bias and does not provide an AGN significance threshold.
These are separate categorical laws, not rebinning of one continuous
dataset. They omit temporal dependence, interpolation, measurement errors,
and fitted quantile boundaries. The deterministic positive controls also
do not represent weak physical signals.

\begin{table}[H]
\centering
\caption{Median joint information $I$ in bits and normalized leakage
$\ell$ under independent categorical variables, from 200 draws per setting.
Exact truth is $I=0$ and $\ell=1$. The 512 tuple setting is an artificial
reference. None of these values is a significance threshold for the AGN.}
\label{tab:cr_calibration}
\small
\input{\SURDResultsPath/calibration_rows.tex}
\end{table}

Some binary positive controls also show unexpected component allocations.
In four preserved examples, reordering predictor histogram axes together
with their physical labels changes the normalized allocation substantially
while leaving total information unchanged. An execution trace of the
unchanged routine records resets to zero of positive specific information
after strict comparison differences around machine precision. We assess
the relevance of this numerical concern to the retained observations next.

\subsection{Predictor ordering audit of the retained observations}

For every retained curve, bin setting, and positive lag, we reconstruct
the original histogram and verify reproduction of the saved components
to absolute tolerance $10^{-10}$. Target dates and tuple counts also
agree. We then evaluate all 24 orders of the four predictor axes, moving
their names with their axes and keeping the target axis fixed. Components
are matched by type and the set of named predictors. There are 5760 order
evaluations across 240 lag and bin configurations. The histogram counts,
target dates, bin boundaries, lag range, and algorithm are unchanged.

An allocation is flagged when its total variation from the original order
exceeds $10^{-8}$, a declared numerical comparison tolerance rather than a
significance threshold. Eleven configurations exceed this tolerance
(Table~\ref{tab:cr_order}). The largest total variation is 0.188 in the
1993 red earlier scenario with three bins. Maximum absolute changes in
joint information and leakage across the audit are
$6.66\times10^{-16}$ bits and $1.83\times10^{-15}$, respectively.
Thus stable total information and component sum normalization do not ensure
stable allocation to named components.

\begin{table}[H]
\centering
\caption{Predictor ordering sensitivity of the retained astronomical
histograms. Each row contains 30 positive lag configurations, each checked
in all 24 predictor orders. Affected lags have component total variation
above $10^{-8}$ in at least one order. The final column is the maximum
over all lags and orders in that row.}
\label{tab:cr_order}
\small
\input{\SURDResultsPath/order_rows.tex}
\end{table}

Both 1993 common curves are stable to predictor ordering within this
tolerance. Their strong bin sensitivity therefore cannot be attributed
solely to this ordering issue. The figures and extrema retain the
original saved ordering as computational diagnostics. We do not select a
preferred permutation or alter the authors' routine to obtain a physical
result. The ordering concern is present in a subset of the observed
configurations and limits physical interpretation of individual allocations.

\subsection{Astrophysical scope and remaining analysis}

The executed curves describe future continuum prediction from its own
history and past line components. They do not provide a supported fixed
date curve of future line response to past continuum in the declared lag
range. A predictive association also requires consideration of a shared
ionizing driver, measurement noise, temporal dependence, and the imperfect
relation between observed optical continuum and illuminating radiation
before it can receive a causal interpretation. Total and component line
measurements retain their shared spectral provenance.

\citet{CRLu2022} report emission line reverberation relative to the optical
continuum in five modern seasons and changes in velocity resolved lag
profiles. \citet{CRXi2025} extend the program to 2023 and report temporal
evolution of the inferred gas kinematics. Those studies motivate analysis
of separate campaigns. Their line response delays and kinematic
interpretations are different quantities from the normalized information
component extrema reported here. Our historical epochs, wavelength
partitions, observed time frame, and retained target direction also differ.
The present diagnostics therefore do not establish agreement or disagreement
with their physical delay measurements or with an inflow or outflow model.

The completed calculations establish substantial sensitivity to bin choice
and retained dates. Neither bin setting is selected as reliable, and no
stable physical delay or causal detection is claimed. No significance
test or flux uncertainty propagation has been performed for these fixed
support curves. The initial categorical calibration and the complete
ordering audit of the retained curves are finished. Calibration matching
the temporal and observational properties remains unperformed. The
expanded predictor construction of thesis
equation~3.18 is implemented and tested, but only its support audit has
been executed; an expanded scientific decomposition is not reported.
These limitations delimit the present results section and the remaining
collaborator analysis. The exploratory analysis is frozen at this stage
for collaborator review. Further scientific runs, any change to the
authors' routine, and any revised probability estimation or lag settings
require a new agreed analysis scope.

\clearpage
\begin{figure}[p]
\centering
\includegraphics[width=.96\textwidth]{\SURDResultsPath/figures/1992_common.png}
\caption{Individual normalized SURD components and information leakage
for the 1992 common scenario. The future target is continuum and predictors
are past continuum, blue, core, and red H$\beta$. Every base lag uses the
same 36 target dates. $F$, $B$, $C$, and $R$ denote those four variables.
The outer $U$, $R$, and $S$ identify the information type. Dotted curves
use two bins and solid curves use three bins. Leakage has its own target
entropy denominator. The curves are descriptive estimates without
significance or physical delay claims.}
\label{fig:cr_1992_common}
\end{figure}

\begin{figure}[p]
\centering
\includegraphics[width=.96\textwidth]{\SURDResultsPath/figures/1993_common.png}
\caption{As in Figure~\ref{fig:cr_1992_common}, for the 1993 common
scenario with 90 fixed continuum target dates. This is the retained curve
with the largest native target support. The changes in individual maxima
between bin settings are summarized in Table~\ref{tab:cr_extrema}.}
\label{fig:cr_1993_common}
\end{figure}

\begin{figure}[p]
\centering
\includegraphics[width=.96\textwidth]{\SURDResultsPath/figures/1993_blue_earlier.png}
\caption{As in Figure~\ref{fig:cr_1992_common}, for the 1993 blue earlier
scenario with 78 fixed continuum target dates. The blue predictor is
evaluated five days earlier than the other predictors. This stipulated
offset is a sensitivity choice, not a measured delay. Its retained dates
differ from those in the common scenario, so scenario differences cannot
be attributed solely to the predictor offset.}
\label{fig:cr_1993_blue}
\end{figure}

\begin{figure}[p]
\centering
\includegraphics[width=.96\textwidth]{\SURDResultsPath/figures/1993_red_earlier.png}
\caption{As in Figure~\ref{fig:cr_1992_common}, for the 1993 red earlier
scenario with 78 fixed continuum target dates. The red predictor is
evaluated five days earlier than the other predictors. This is a
stipulated sensitivity choice. Dates are held fixed within this curve,
but are not forced to match the other scenarios.}
\label{fig:cr_1993_red}
\end{figure}
\clearpage
